%% annex/p07/lambda-calculi.tex --- proof material for P7 chapter 7.4
%% New content: no source-compression denominator applies.

\begin{annexitem}{annex:p7.lambda.1}{thm:p07-church-rosser}{Church--Rosser by parallel reduction}
\begin{proof}
All terms in Method~\ref{def:p07-parallel-reduction} are read modulo
alpha-equivalence.  First observe, by induction on a term, that
\(M\mathrel{\mathsf{Par}}M\): use the variable clause at a variable and the
application and abstraction clauses at the two constructors.  The same
induction shows that complete development respects alpha-equivalence.  The
only nonliteral step is an abstraction: after replacing
\(\lambda y.P\) by \(\lambda z.P[z/y]\), with \(z\) fresh, the induction
hypothesis gives
\(\operatorname{dev}(P[z/y])=_{\alpha}
  \operatorname{dev}(P)[z/y]\).  Whether an application has an abstraction at
its head is unchanged by alpha-renaming, so the two application clauses of
\(\operatorname{dev}\) agree on alpha-equivalence classes.

We next prove the parallel-substitution lemma:
\[
 M\mathrel{\mathsf{Par}}M',\quad
 N\mathrel{\mathsf{Par}}N'
 \quad\Longrightarrow\quad
 M[N/x]\mathrel{\mathsf{Par}}M'[N'/x].                 \tag{1}
\]
Before the induction, alpha-rename every binder in the displayed parallel
derivation away from \(x\) and from the free variables of \(N,N'\); the first
paragraph shows that this changes neither hypothesis nor conclusion.  Induct on the displayed derivation of
\(M\mathrel{\mathsf{Par}}M'\).  If \(M=x\), both sides of (1) are
\(N\) and \(N'\), so use the second hypothesis.  If \(M=y\ne x\), both
sides are \(y\), so use reflexivity.  The application case follows by applying
the induction hypotheses to the two subterms and then the application clause.
For an abstraction \(\lambda y.P\), freshness gives \(y\ne x\); apply the
induction hypothesis to \(P\) and use the abstraction clause.  In the beta
case write the source as \((\lambda y.P)Q\) and the target as
\(P'[Q'/y]\).  The induction hypotheses give
\(P[N/x]\mathrel{\mathsf{Par}}P'[N'/x]\) for the body and the analogous
statement for the argument.  If the redex binder is \(y\), the beta clause
then has target
\(P'[N'/x][Q'[N'/x]/y]\), which is alpha-equal to
\((P'[Q'/y])[N'/x]\) because \(y\notin\FV(N')\).  This
composition-of-substitutions identity follows by induction on \(P'\),
alpha-renaming an abstraction binder fresh for all substituted terms; thus no
textual substitution or capture is used.

The development lemma is
\[
 M\mathrel{\mathsf{Par}}P
 \quad\Longrightarrow\quad
 P\mathrel{\mathsf{Par}}\operatorname{dev}(M).          \tag{2}
\]
Induct on the parallel derivation.  The variable case is immediate.  The
abstraction case follows from the induction hypothesis and the abstraction
clause.  For an application-congruence conclusion
\(MN\mathrel{\mathsf{Par}}M'N'\), use the two induction hypotheses.  If
\(M\) is not an abstraction, the application clause gives
\(M'N'\mathrel{\mathsf{Par}}
  \operatorname{dev}(M)\operatorname{dev}(N)\), which is
\(\operatorname{dev}(MN)\).  If \(M=\lambda x.Q\), the induction hypotheses
give
\(M'=_{\alpha}\lambda x.Q'\) with
\(Q'\mathrel{\mathsf{Par}}\operatorname{dev}(Q)\), and
\(N'\mathrel{\mathsf{Par}}\operatorname{dev}(N)\); the beta clause therefore
gives
\[
 M'N'\mathrel{\mathsf{Par}}
 \operatorname{dev}(Q)[\operatorname{dev}(N)/x]
 =\operatorname{dev}(MN).
\]
Finally, if the last rule is the parallel beta clause, its conclusion is
\(Q'[N'/x]\).  The induction hypotheses for the body and argument, followed
by (1), give
\[
 Q'[N'/x]\mathrel{\mathsf{Par}}
 \operatorname{dev}(Q)[\operatorname{dev}(N)/x],
\]
again exactly the complete development of the source redex.  This exhausts
the four clauses and proves (2).

If
\(M\mathrel{\mathsf{Par}}N_0\) and
\(M\mathrel{\mathsf{Par}}N_1\), (2) supplies
\(N_i\mathrel{\mathsf{Par}}\operatorname{dev}(M)\) for both \(i\).
Thus one-step parallel reduction has the diamond property.

It remains to relate the auxiliary relation to ordinary beta reduction.  A
root beta contraction is a parallel beta step by reflexivity of the developed
body and argument; induction on a term context closes this fact under
application and abstraction.  Hence
\(\longrightarrow_\beta\subseteq\mathsf{Par}\).  Conversely, induction on a
parallel derivation sequentializes it.  The variable case takes no beta step;
the application and abstraction cases concatenate the finite reductions of
their subterms inside the corresponding contexts; and the parallel beta case
first reduces the body and argument inside
\((\lambda x.M)N\) to \((\lambda x.M')N'\) and then contracts the root redex
to \(M'[N'/x]\).  Therefore
\(\mathsf{Par}\subseteq\longrightarrow_\beta^{*}\), and the two relations
have the same reflexive--transitive closure.

For completeness, lift the one-step diamond to finite paths explicitly.  Given
paths of lengths \(m,n\) from one term, place them along the left and top edges
of an \(m\)-by-\(n\) grid.  Induct lexicographically on the cell coordinates.
At each cell the already constructed top and left edges are two parallel steps
with a common source; the diamond property supplies the bottom and right edges.
After finitely many cells, the lower-right corner is reachable from both path
endpoints.  Replacing each parallel edge by its finite beta simulation gives a
common beta reduct.  This is exactly the assertion of
Theorem~\ref{thm:p07-church-rosser}.
\end{proof}
\end{annexitem}

\begin{annexitem}{annex:p7.lambda.2}{thm:p07-beta-normal-form-unique}{Uniqueness of beta-normal forms}
\begin{proof}
Suppose
\(M\longrightarrow_{\beta}^{*}N_1\) and
\(M\longrightarrow_{\beta}^{*}N_2\), with both \(N_1,N_2\) beta-normal.
Theorem~\ref{thm:p07-church-rosser} gives \(P\) such that
\(N_i\longrightarrow_{\beta}^{*}P\) for \(i=1,2\).  No nonreflexive beta step
can start from a beta-normal term, so each of these two reductions consists
only of the reflexive step, modulo the bound-variable representatives fixed
by Definition~\ref{def:p07-alpha-substitution}.  Thus \(N_1\) and \(N_2\) are
alpha-equivalent.
\end{proof}
\end{annexitem}
